\section{Related Work}

\paragraph{Scientific peer review and review assistance.}
Peer-review datasets have enabled computational study of review text, scores, decisions, and reviewing workflows.
PeerRead introduced a public corpus of paper drafts, decisions, and expert reviews from venues including ACL, NeurIPS, and ICLR \citep{kang-etal-2018-dataset}.
NLPeer expands this line with a clearly licensed, multi-domain resource that includes drafts, camera-ready versions, reviews, metadata, and versioning information \citep{dycke-etal-2023-nlpeer}.
More recent resources such as Re2 collect full-stage review and rebuttal discussions across many OpenReview venues \citep{zhang-etal-2025-re2}.
Task inventories for peer-review automation also emphasize that reviewing decomposes into many different subtasks and that dataset documentation matters for comparability \citep{staudinger-etal-2024-analysis}.
Earlier automatic-review systems such as DeepReviewer and ReviewRobot also establish the paper-review automation line with score/comment generation objectives \citep{leng-etal-2019-deepreviewer,wang-etal-2020-reviewrobot}.
Recent LLM-era work extends this line with larger review datasets and staged reasoning pipelines for review generation \citep{zhu-etal-2025-deepreview}.
These resources provide the raw setting for studying scientific review, but their main tasks are usually paper-level prediction, review understanding, or discussion modeling.
RevTrack changes the unit of evaluation.
Rather than predicting acceptance, summarizing reviews, or modeling discussions wholesale, it isolates one reviewer concern and asks whether later revision evidence changes that concern's status.
This makes revision tracking closer to an issue ledger than to a static review corpus.

\paragraph{Review-question answering and evidence grounding.}
PeerQA uses peer-review questions as document-level scientific QA prompts and asks systems to retrieve evidence, identify unanswerable questions, and generate answers \citep{baumgartner-etal-2025-peerqa}.
RevTrack shares the emphasis on authentic reviewer information needs, but changes the output from answer generation to status tracking.
More broadly, evidence-verification benchmarks such as FEVER test whether claims are supported, refuted, or lack enough information with respect to textual evidence \citep{thorne-etal-2018-fever}.
RevTrack differs in two ways.
First, the claim is not a standalone fact: it is an old review concern whose status depends on later author response and revision evidence.
Second, the correct label is not only supported versus refuted.
Partial fixes and regressions are first-class outcomes because scientific revision often narrows, complicates, or worsens a concern rather than simply resolving it.
This perspective is also connected to QA with conflicting evidence, where systems must avoid over-confident answers when contexts disagree \citep{liu-etal-2025-open,shaier-etal-2024-adaptive}.

\paragraph{LLM-assisted reviewing.}
LLM-based review assistance has focused heavily on generating reviewer feedback.
Large-scale empirical work suggests that GPT-4 feedback can overlap with human-reviewer feedback and can be perceived as useful by authors \citep{liang-etal-2023-can}.
Targeted critique/meta-review assistance has also been studied directly for NLP papers \citep{du-etal-2024-llms-assist}.
At the same time, pilot evidence reports only limited net utility for generic peer-review assistance prompts \citep{robertson-2023-gpt4-helpful}.
For example, MARG uses multiple specialized LLM agents to generate scientific-paper feedback and reduce generic comments \citep{darcy-etal-2024-marg}.
This direction is useful, but review generation alone does not test whether an assistant can retire stale criticisms after a paper changes.
RevTrack therefore evaluates a narrower and more auditable capability: revision-aware issue tracking.
This framing also separates assistance quality from fluency.
A system can write a plausible comment while still failing the longitudinal decision of whether a prior comment should remain active.

\paragraph{Reliability, calibration, and hallucination under evidence conflict.}
LLM-as-judge methods have become common for scalable open-ended evaluation \citep{zheng-etal-2023-judging,liu-etal-2023-g}, and related work explores using panels of smaller judges to reduce cost and single-model bias \citep{verga-etal-2024-replacing}.
At the same time, LLM evaluators can show position and fairness biases that require calibration \citep{wang-etal-2024-large-language-models-fair}.
LLMs are also increasingly studied as data annotators or annotation assistants \citep{tan-etal-2024-large,gu-etal-2025-large}.
These lines motivate our prompted-LLM baselines, but they also make direct substitution of LLM labels for validated status labels risky.
Recent work has shown that strong language models can remain unreliable under factual pressure or conflicting evidence, even when outputs are fluent \citep{lin-etal-2022-truthfulqa,ji-etal-2023-hallucination-survey}.
Confidence-aware analyses also show that model certainty is not a sufficient proxy for correctness \citep{kadavath-etal-2022-language}.
This reliability gap under shift is consistent with broader calibration and robustness findings beyond review tasks \citep{guo-etal-2017-calibration,ovadia-etal-2019-shift,koh-etal-2021-wilds}.
Post-hoc hallucination checks can improve auditing \citep{manakul-etal-2023-selfcheckgpt}, but they do not directly solve the longitudinal decision in review workflows.
RevTrack complements this line by making the evaluation target explicit: whether a previously raised concern should stay active after revision evidence.

\paragraph{Evaluation under skew.}
The accuracy trap observed here is a standard risk in imbalanced classification, but it is especially damaging for review support.
The minority labels are often the scientifically important ones: fixed cases determine whether stale criticism should be retired, and unresolved cases determine whether a blocker remains active.
RevTrack therefore treats macro-F1, fixed-case recovery, unresolved-case recovery, and regressed-case analysis as part of the task definition rather than secondary diagnostics.
